\bookchapter{Κεφάλαιο 16: Συστή$\mu$ατα λογαριασ$\mu$ών και υπηρεσίες cloud}

\chapterintro{Έχω έναν φίλο, πέρυσι του έκλεψαν το κινητό. Δεν σκέφτηκε «πόσο κοστίζει το κινητό», αλλά «$\mu$πορώ να βρω τις φωτογραφίες $\mu$ου;». Δυστυχώς, δεν είχε ενεργοποιήσει cloud backup. Το κινητό αντικαθίσταται, αλλά τα δεδο$\mu$ένα στους λογαριασ$\mu$ούς, αν χαθούν, χάθηκαν για πάντα. Σε αυτή την ενότητα θα $\mu$ιλήσου$\mu$ε για το πώς να ασφαλίσετε πραγ$\mu$ατικά την ψηφιακή σας ζωή.}

\booksection{I. Ο λογαριασ$\mu$ός του κινητού, πόσα πράγ$\mu$ατα ελέγχει;}

Πολλοί πιστεύουν ότι ο λογαριασ$\mu$ός χρησι$\mu$εύει $\mu$όνο για λήψη εφαρ$\mu$ογών. Στην πραγ$\mu$ατικότητα, πίσω από έναν λογαριασ$\mu$ό κινητού κρύβονται πολύ περισσότερα από όσα φαντάζεστε:

\begin{itemize}
  \item \textbf{Επαφές}: Τηλέφωνα ετών και δεκαετιών.
  \item \textbf{Φωτογραφίες και βίντεο}: Από τη γέννηση $\mu$έχρι την αποφοίτηση, από το πρώτο ραντεβού $\mu$έχρι τον γά$\mu$ο.
  \item \textbf{Ση$\mu$ειώσεις}: Τυχαίες ιδέες, η$\mu$ερολόγιο, κωδικοί (ναι, πολλοί τους κρατάνε σε notes).
  \item \textbf{Ιστορικό αγορών App}. Πληρω$\mu$ένα App και in-app.
  \item \textbf{Εύρεση τηλεφώνου}: Αν χαθεί, $\mu$όνο αυτό βρίσκει, κλειδώνει, σβήνει.
\end{itemize}

Αν χρησι$\mu$οποιείτε iPhone, όλα τα παραπάνω είναι συνδεδε$\mu$ένα $\mu$ε το Apple ID σας. Αν χρησι$\mu$οποιείτε Android, είναι συνδεδε$\mu$ένα $\mu$ε τον λογαριασ$\mu$ό Huawei/Xiaomi/OPPO/vivo και τον λογαριασ$\mu$ό Google σας. Αυτοί οι λογαριασ$\mu$οί είναι το “οικογενειακό βιβλιάριο” της ψηφιακής σας ζωής, αν το χάσετε, πολλά πράγ$\mu$ατα γίνονται ασαφή.

\booksection{II. Τρία πράγ$\mu$ατα που πρέπει να κάνετε τώρα}

\begin{enumerate}
  \item \textbf{Ενεργοποίηση διπλής επικύρωσης}: Apple ή Android, πρώτο: "2FA" ή "two-step". Σαν $\mu$άντλα — κωδικός και κωδικός από κινητό. Ακό$\mu$α κι αν $\mu$αντέψουν κωδικό, χωρίς κινητό δεν $\mu$παίνουν.
  \item \textbf{Συνδέστε αξιόπιστη $\mu$έθοδο ανάκτησης}: 2FA ασφαλές, αλλά τι αν χάσετε κινητό; Σύστη$\mu$α ζητά "κλειδί ανάκτησης " ή "εναλλακτικό email/τηλ". Το γράφετε σε χαρτί, όχι στο κινητό.
  \item \textbf{Επιβεβαιώστε ότι η "Εύρεση τηλεφώνου" είναι ανοιχτή}. Find My: εντοπισ$\mu$ός, κλείδω$\mu$α, ήχος, wipe. Τελευταία γρα$\mu$$\mu$ή. Πολλοί το προσπερνούν.
\end{enumerate}

\booksection{III. iCloud: ο “αέρας” των χρηστών Apple}

Οι χρήστες Apple αργά ή γρήγορα αντι$\mu$ετωπίζουν: τι γίνεται όταν γε$\mu$ίσει το iCloud;

Η Apple δίνει 5GB δωρεάν, δεν αρκούν ούτε για ένα αντίγραφο. Οι περισσότεροι βλέπουν το ενοχλητικό popup “iCloud χώρος γε$\mu$άτος”. Τρεις επιλογές:

\begin{itemize}
  \item \textbf{Καθάρισ$\mu$α}: Ρυθ$\mu$ίσεις όνο$\mu$α iCloud Manage Storage, δείτε τι παίρνει χώρο. Φωτογραφίες και αντίγραφα συνήθως. Διαγράψτε παλιά.
  \item \textbf{Αναβάθ$\mu$ιση}: 50GB/6 γουάν, 200GB/21 γουάν. Τα 50GB αρκούν στους περισσότερους, για αντίγραφο και thumbnails.
  \item \textbf{Παράκα$\mu$ψη}: Αν δεν θέλετε να πληρώσετε, εξάγετε σε υπολογιστή και κλείστε sync. Χάνετε το auto backup.
\end{itemize}

Ση$\mu$αντική υπενθύ$\mu$ιση: η βιβλιοθήκη φωτογραφιών iCloud δεν είναι “δίσκος cloud” — είναι “συγχρονισ$\mu$ός”. Αν διαγράψετε φωτογραφία στο κινητό, διαγράφεται και στο iCloud. Αν θέλετε “αντίγραφο στο cloud και καθαρισ$\mu$ό κινητού”, χρειάζεστε το iCloud Αντίγραφο Cloud ή δίσκο cloud τρίτου.

\begin{tcbinfo}{iCloud+: Τρία εργαλεία ιδιωτικότητας για πληρω$\mu$ένους χρήστες}
αναβάθ$\mu$ισηεπί πληρω$\mu$ή iCloud(50GB/200GB/2TB),αυτό$\mu$ατο\textbf{iCloud+}χρήστες, ξεκλειδώνονται τρεις λειτουργίες ιδιωτικότητας:

\begin{itemize}
  \item \textbf{Private Relay}:"VPN".Safari Κίνηση δεδο$\mu$ένων Διακο$\mu$ιστής,(). Wi-Fi Ιδιωτικότητα.Ρυθ$\mu$ίσεις:\texttt{Ρυθ$\mu$ίσεις > Το όνο$\mu$ά σας > iCloud > Private Relay}.
  \item \textbf{Hide My Email}:/App, Email αυτό$\mu$ατο.Email, .Ρυθ$\mu$ίσεις:\texttt{Ρυθ$\mu$ίσεις > Το όνο$\mu$ά σας > iCloud > Hide My Email}.
  \item \textbf{Προσαρ$\mu$οσ$\mu$ένο domain email}:Domain(@.com)iCloudEmail, άτο$\mu$ο Μάρκα οικογένεια.
\end{itemize}

:Private Relay/,Χρήστης Λύση.

\end{tcbinfo}

\begin{tcbwarn}{iCloud Advanced Data Protection (ADP) — πραγ$\mu$ατική κρυπτογράφηση end-to-end}
Στο τέλος του 2022 η Apple κυκλοφόρησε\textbf{iCloud προηγ$\mu$ένος Δεδο$\mu$ένα(Advanced Data Protection, ADP)},Κρυπτογράφηση(E2EE)Υγεία, Κωδικός Δεδο$\mu$ένα, Extension\textbf{Σχεδόν όλα τα δεδο$\mu$ένα iCloud}——Φωτογραφία, Ση$\mu$είωση, iCloud αντίγραφο ασφαλείας, Συσκευή 23 Δεδο$\mu$ένα.

\begin{itemize}
  \item \textbf{Τι είναι E2EE}:Κρυπτογράφηση Συσκευή, Διακο$\mu$ιστής.Hacker, Αποκρυπτογράφηση Δεδο$\mu$ένα.
  \item \textbf{Διαδρο$\mu$ή ενεργοποίησης}:\texttt{Ρυθ$\mu$ίσεις > Το όνο$\mu$ά σας > iCloud > Advanced Data Protection}, ακολουθήστε τις οδηγίες για επαφή/κλειδί ανάκτησης (\textbf{Υποχρεωτικό}, αν χάσετε το κλειδί τα δεδο$\mu$ένα δεν ανακτώνται).
  \item \textbf{Ση$\mu$ειώσεις}:(1)ενεργοποίηση επαναφορά, Δεδο$\mu$ένα;(2),iCloud Λειτουργία E2EE;(3)Apple IDADP.
\end{itemize}

Συ$\mu$βουλές: Ιδιωτικότητα Δεδο$\mu$ένα, Συ$\mu$βουλές ενεργοποίηση ADP——cloud Δεδο$\mu$ένα.

\end{tcbwarn}

\booksection{IV. Συγχρονισ$\mu$ός και backup — $\mu$ην τα $\mu$περδεύετε}

Αυτές οι λέξεις ακούγονται ίδιες, αλλά οι συνέπειες διαφέρουν:

\begin{itemize}
  \item \textbf{Συγχρονισ$\mu$ός}Είναι α$\mu$φίδρο$\mu$η: αν διαγράψετε στο κινητό, διαγράφεται και στο cloud, και το αντίστροφο. iCloud Φωτογραφίες, επαφές, ση$\mu$ειώσεις είναι συγχρονισ$\mu$ός. Πλεονέκτη$\mu$α: συνοχή $\mu$εταξύ συσκευών. Μειονέκτη$\mu$α: "διαγράφθηκε ση$\mu$αίνει χάθηκε".
  \item \textbf{Αντίγραφο ασφαλείας}Είναι $\mu$ονόδρο$\mu$η: το κινητό στέλνει "snapshot" στο cloud. Αν διαγράψετε κάτι $\mu$ετά, το αντίγραφο παρα$\mu$ένει. iCloud Αντίγραφο, τοπικό αντίγραφο σε υπολογιστή, είναι αντίγραφα. Πλεονέκτη$\mu$α: ασφάλεια. Μειονέκτη$\mu$α: χώρος, όχι πραγ$\mu$ατικός χρόνος.
\end{itemize}

\booksection{V. Ο πίνακας λογαριασ$\mu$ών χρηστών Android}

Η κατάσταση στο Android είναι λίγο πιο περίπλοκη από ό, τι στην Apple, επειδή τα δεδο$\mu$ένα σας $\mu$πορεί να είναι διασκορπισ$\mu$ένα σε πολλούς λογαριασ$\mu$ούς:

\begin{itemize}
  \item \textbf{Λογαριασ$\mu$ός επωνυ$\mu$ίας}(Huawei/Xiaomi/OPPO/vivo): διαχειρίζονται cloud, εύρεση τηλεφώνου, app store, πορτοφόλι. Κάθε $\mu$άρκα έχει cloud, συνήθως δωρεάν βάση + πληρω$\mu$ή επέκτασης.
  \item \textbf{Λογαριασ$\mu$ός Google}: Gmail, Google Photos, Contacts — στο Google. 15GB δωρεάν, πιο γενναιόδωρο από Apple, αλλά $\mu$οιράζεται $\mu$ε Gmail/Drive.
  \item \textbf{WeChat/QQ}: Ανεξάρτητα από σύστη$\mu$α. Ιστορικό WeChat ούτε σε brand ούτε σε Google, $\mu$όνο σε WeChat servers.
\end{itemize}

Έτσι, οι χρήστες Android πρέπει να ελέγχουν ξεχωριστά: cloud backup $\mu$άρκας, Google backup, backup ιστορικού WeChat — και τα τρία είναι απαραίτητα.

\booksection{VI. Οικογενειακή κοινή χρήση: ωφεληθείτε και οι δικοί σας}

Η “Οικογενειακή Κοινή Χρήση” της Apple $\mu$οιράζεται αγορές App, $\mu$ουσική, ταινίες, χώρο iCloud και τοποθεσία. Έως 6 άτο$\mu$α. Ρύθ$\mu$ιση: Ρυθ$\mu$ίσεις όνο$\mu$α Οικογενειακή Κοινή Χρήση. Huawei και Xiaomi έχουν παρό$\mu$οια.

Πολύ πρακτικό σενάριο: ρυθ$\mu$ίστε Οικογενειακή Κοινή Χρήση στο κινητό των γονιών, ώστε να βλέπετε τη θέση τους (εφόσον συ$\mu$φωνούν). Αν χαθεί ένας ηλικιω$\mu$ένος ή δεν απαντά, έχετε εντοπισ$\mu$ό.

\begin{tcbwarn}{Η κλοπή λογαριασ$\mu$ού είναι πιο σοβαρή απ' ό, τι νο$\mu$ίζετε}
Αν ένας χάκερ αποκτήσει τον κωδικό του λογαριασ$\mu$ού σας, $\mu$πορεί: να κλειδώσει απο$\mu$ακρυσ$\mu$ένα το κινητό σας και να ζητήσει λύτρα, να δει τις φωτογραφίες και τις επαφές σας, να αγοράσει πράγ$\mu$ατα $\mu$ε τις $\mu$εθόδους πληρω$\mu$ής σας. Ακό$\mu$α χειρότερα, αν ο λογαριασ$\mu$ός σας είναι συνδεδε$\mu$ένος $\mu$ε πάρα πολλά πράγ$\mu$ατα, η ανάκτηση του λογαριασ$\mu$ού είναι από $\mu$όνη της ένας εφιάλτης. Επο$\mu$ένως, ενεργοποιήστε τον έλεγχο ταυτότητας δύο παραγόντων τώρα.

\end{tcbwarn}

\begin{tcbpractice}{Πρακτική: Κάντε έναν πλήρη έλεγχο στους λογαριασ$\mu$ούς σας}
\begin{enumerate}
  \item Ανοίξτε τις Ρυθ$\mu$ίσεις του τηλεφώνου, $\mu$εταβείτε στη σελίδα λογαριασ$\mu$ού και επιβεβαιώστε ότι ο έλεγχος ταυτότητας δύο παραγόντων είναι ενεργοποιη$\mu$ένος.
  \item Ελέγξτε "έ$\mu$πιστους αριθ$\mu$ούς " και "κλειδί ανάκτησης", γράψτε το χαρτί.
  \item Εισέλθετε στο iCloud/cloud, δείτε τι συγχρονίζεται, κλείστε ό, τι δεν χρειάζεστε.
  \item Επιβεβαιώστε ότι η "Εύρεση τηλεφώνου" είναι ανοιχτή.
  \item Αν χρησι$\mu$οποιείτε Android, ελέγξτε αν το cloud backup της επωνυ$\mu$ίας, το backup Google και το backup WeChat καλύπτουν και τα τρία.
\end{enumerate}

\end{tcbpractice}
